\documentclass[10pt,a4paper]{amsart}
\usepackage[margin=19mm]{geometry}
\usepackage{amsmath,amssymb,mathtools}
\usepackage[scr=rsfs,cal=euler]{mathalfa}
\usepackage{xfrac}
\usepackage[unicode,hidelinks,bookmarks=false]{hyperref}
\newcommand{\ie}{i.e.\ }
\newcommand{\cf}{cf.\ }
\newcommand{\resp}{respectively\ }
\newcommand{\Subsection}{\subsection}
\providecommand{\nobreakdash}{\nobreak\hbox{-}}
\setlength{\emergencystretch}{2em}
\multlinegap0pt
\usepackage{amssymb}
\usepackage{enumerate}
\newtheorem{theorem}{Theorem}
\newtheorem{lemma}{Lemma}
\newtheorem{proposition}{Proposition}
\title{On Extensions of $D(4)$-triples by Adjoining Smaller Elements}
\author{Marija Bliznac Trebje\v{s}anin, Pavao Radi\'c}
\date{Source: arXiv:2306.00850v2 (CC BY 4.0)}
\begin{document}
\maketitle

\noindent\emph{Source and licence.} On Extensions of $D(4)$-triples by Adjoining Smaller Elements. Version arXiv:2306.00850v2 (CC BY 4.0), distributed by arXiv under the Creative Commons Attribution 4.0 International licence, http://creativecommons.org/licenses/by/4.0/, as recorded in the arXiv licence metadata for this exact version and shown on the abstract page https://arxiv.org/abs/2306.00850v2. Full licence notice: \texttt{licenses/arxiv-2306.00850-CC-BY-4.0.txt}.\par\smallskip

\noindent\textit{Editorial scope.} Counted unit: the article's theorem theorem (source0.tex lines 574--578). The article's complete original proof of it is delivered with it (source0.tex lines 580--644, one \texttt{proof} environment of 65 lines). No mathematics is written, repaired or re-derived: the statement and the proof are the article's own lines, in the article's own notation, and no retained line is substituted or re-flowed.  Retained uncounted as the source-local context the proof needs: lines 225--230; lines 241--244; lines 246--249; lines 271--273; lines 553--556.  Nothing was omitted from any retained span, so the package carries no omission record.  No retained line contains a citation, so the wrapper carries no bibliography and no bibliography entry.  No retained line contains a figure environment, an image-inclusion command or drawing code, so no image is delivered and the package declares no asset dependency.  Editorial formatting only, in addition: an AMS \texttt{amsart} preamble that restates the article's own declarations and macros used by the retained lines, namely the environments \texttt{theorem}, \texttt{lemma}, \texttt{proposition} on separate counters, together with the standard packages those lines need.  The retained environments print in this document as Theorem 1, Lemma 1, Proposition 1 in order of appearance, whereas the article prints the same items with the numbering of its own full document.  The complete native source of the article as distributed in the arXiv e-print for this exact version is bundled unchanged as \texttt{sources/142375/source0.tex}.
\begin{theorem}\label{tm: gornje_ograde}
Let $a_1,a_2$ be integers such that $0<a_1 \leq a_2-2,\ a_2 \geq 3$ and let's denote $a_1' := \max\{4(a_2-a_1),4a_1\}$. Also, let $N$ be a multiple of $a_1a_2$ and assume that $N \geq 396.59a_1'a_2^2(a_2-a_1)^2$. Then the numbers $\theta_1 = \sqrt{1+\frac{4a_2}{N}}$ and $\theta_2 = \sqrt{1+\frac{4a_1}{N}}$ satisfy
$$ \left\{ \left| \theta_1-\frac{p_1}{q} \right| , \left|\theta_1-\frac{p_2}{q} \right|     \right\} > \frac{a_1}{3.53081 \cdot 10^{27} a_1'a_2N}q^{-\lambda}$$
for all integers $p_1,p_2,q$, with $q>0$, where 
\[\lambda=1+\frac{\log(2.500788a_1^{-1}a_1'a_2N)}{\log(0.04216a_1^{-1}a_2^{-1}(a_2-a_1)^{-2}N^2)}<2.\]
\end{theorem}

\begin{lemma}\label{gornje_ograde3}
     Assume that $c >4.1\cdot10^{-5}a_2b^3$ and $b>10^5$. If $z=w_n$ with $n \geq 4$, then 
    \[ \log z > \frac{n}{2}\log(bc) .\]
\end{lemma}  

\begin{lemma}\label{lemma: gornja_n_1}
    Let $\{a_1,b,c,d\}$ and $\{a_2,b,c,d\}$ be $D(4)$-quadruples such that the assumptions in Theorem \ref{tm: gornje_ograde} and Lemma \ref{gornje_ograde3} hold. If $z=v_m=w_n$ has a solution for some integers $m$ and $n$ with $n \geq 4$, then 
   $$ n < \frac{8\log(8.4034 \cdot 10^{13} a_1^{1/2}(a_1')^{1/2}a_2^2c)\log(0.20533a_1^{1/2}a_2^{1/2}(a_2-a_1)^{-1}c)}{\log(bc)\log(0.01685a_1(a_1')^{-1}a_2^{-1}(a_2-a_1)^{-2}c)}. $$
\end{lemma}

\begin{lemma}\label{lema: donja na n u terminima b i c}
    Assume that $c>\max\{0.56b^3,0.001a_2^2b^2\}$ and $b>10^5$. If $z=v_m=w_n$ has a solution for some integers $m$ and $n$ with $n\geq 4$, then $m\equiv n\ (\bmod \ 2)$ and $n>0.09226b^{-1/2}c^{1/4}$ for odd $n$ and $n>0.340134b^{-1/2}c^{1/2}$ for even $n$.
\end{lemma}

\begin{proposition} 
\label{prop: prop1}
    Assume that $\{a_1,b,c,d\}$ and $\{a_2,b,c,d\}$ are $D(4)$-quadruples with $a_1<a_2<b<c<d$. Then $a_2>4a_1$.
\end{proposition}

\begin{theorem}
    Assume that $\{a_1,b,c,d\}$ and $\{a_2,b,c,d\}$ are $D(4)$-quadruples with $a_1<a_2<b<c<d$. Then 
       $$a_2>0.0251a_1^2.$$
   
\end{theorem}

\begin{proof}    
We will first prove that $a_2>0.01a_1^2$. 
Assume to the contrary, that $a_2\leq 0.01a_1^2$. From Proposition \ref{prop: prop1} we have $a_2>4a_1$. This implies $a_1\geq 401$ and $a_2\geq 1605.$ 
It is easy to see that the assumptions of Lemmas \ref{lemma: gornja_n_1} and \ref{lema: donja na n u terminima b i c} are now satisfied.

 Notice that $3a_1<a_2-a_1<a_2\leq0.01a_1^2$. We obtain the following inequalities: 
    \begin{align*}
 8.4034\cdot 10^{13}a_1^{1/2}(a_1')^{1/2}a_2^2c&<1.6807\cdot 10^9a_1^{11/2}c\\
 0.20533a_1^{1/2}a_2^{1/2}(a_2-a_1)^{-1}c&<0.00685a_1^{1/2}c\\
 0.01685a_1(a_1')^{-1}a_2^{-1}(a_2-a_1)^{-2}c&>421250a_1^{-7}c,
\end{align*}
so we observe an inequality
\begin{equation}\label{eq:njdn_manja001}
    n<\frac{8\log(1.6807\cdot 10^{9}\cdot a_1^{11/2}c)\log(0.00685a_1^{1/2}c)}{\log(bc)\log(421250a_1^{-7}c)}.
\end{equation}

We will separate the proof in cases depending on relations between $b$ and $a_1^2$. 

\textbf{Case 1:} $b \geq a_1^2$. 

Here we have $c>0.25a_1^8$. The right-hand side of inequality (\ref{eq:njdn_manja001}) is decreasing in $c$ for $c>421250^{-1}a_1^7$ which obviously holds. Also, $n>0.09226\cdot 0.25^{1/4}a_1$. Applying all those inequalities to (\ref{eq:njdn_manja001}) yields $a_1\leq 552$. 

Now we can rewrite the inequality (\ref{eq:njdn_manja001}) and together with  $n>0.2465c^{1/12}$ 
obtain the numerical upper bound for  $c$, i.e. $c<2.2666\cdot 10^{24}$. This implies $b<2.086\cdot 10^{8}\cdot a_1^{-2/3}$. We consider these finitely many pairs $\{a_1,a_2\}$, $a_2\leq 0.01a_1^2$ and a computer search for extensions within these bounds for $b$ and $c$ returns that there are no such numbers.

\textbf{Case 2:} $b<a_1^2$. 

Observe that, since $a_1<a_2/4<b/4$, $\frac{3}{4}a_2<a_2-a_1<a_2\leq 0.01a_1^2$ and $0.25b^4<c<0.25b^5$ we have
   \begin{align*}
 8.4034\cdot 10^{13}a_1^{1/2}(a_1')^{1/2}a_2^2c&<8.4034\cdot 10^{13}b^3c<2.101\cdot 10^{13}b^8\\
 0.20533a_1^{1/2}a_2^{1/2}(a_2-a_1)^{-1}c&<0.137c<0.03425b^5\\
 0.01685a_1(a_1')^{-1}a_2^{-1}(a_2-a_1)^{-2}c&>1.725\cdot 10^7b^{-7}c>4.3125\cdot 10^6b^{-3}. 
\end{align*}

Also, $n>0.09226b^{-1/2}c^{1/4}>0.09226\cdot 0.25^{1/4}b^{1/2}$. Applying this to inequalities from Lemmas \ref{lemma: gornja_n_1} and \ref{lema: donja na n u terminima b i c} yields $b\leq 162$ which is a contradiction. Thus, we have proved that $a_2>0.01a_1^2$.


Assume now that $a_2 \leq 0.0251a_1^2$. From Proposition \ref{prop: prop1} we have $a_2>4a_1$ and get $a_1\geq 160$ and $a_2 > 4\cdot 160=640$. It is easy to see that the assumptions of Lemmas \ref{lemma: gornja_n_1} and \ref{lema: donja na n u terminima b i c} are satisfied.

We follow the steps from the first part of the proof.
    Notice that $0.003a_1^2<a_2-a_1<a_2\leq0.0251a_1^2$ holds. 
so we obtain an inequality
\begin{equation}\label{eq:njdn_manja00251}
    n<\frac{8\log(1.678\cdot 10^{10}a_1^{11/2}c)\log(0.3286a_1^{-1/2}c)}{\log(bc)\log(10613a_1^{-7}c)}.
\end{equation}

As before, we consider two cases.

\textbf{Case 1:} $b \geq a_1^2$. Since the same argument holds as before, 
 we can use $c>10613^{-1}a_1^7$ and $n>0.09226\cdot 0.25^{1/4}a_1$. Applying all those inequalities yields $a_1\leq 682$. 

Now we can rewrite the inequality (\ref{eq:njdn_manja00251}) by inserting numerical bounds for $a_1$ and together with $n>0.2465c^{1/12}$ obtain $c<3.06\cdot 10^{24}$ and $b<2.205\cdot 10^{8}\cdot a_1^{-2/3}$. We consider the finitely many pairs $\{a_1,a_2\}$ which satisfy $ 0.01a_1^2<a_2\leq 0.0251a_1^2$ and search for extensions within the bounds for $b$ and $c$. A computer search as described before returns that there is no $c$ within the given bounds.

\textbf{Case 2:} $b<a_1^2$. 
Because of $b>10^5$, this assumption yields $a_1 \geq 317$. 


Observe that, since $a_1<a_2/4<b/4$, $0.007b<0.007a_1^2<a_2-a_1\leq 0.0251a_1^2$ for $a_1\geq 317$, and $0.25b^4<c<0.25b^5$ we have
   \begin{align*}
 8.4034\cdot 10^{13}a_1^{1/2}(a_1')^{1/2}a_2^2c&<8.4034\cdot 10^{13}b^3c<2.101\cdot 10^{13}b^8\\
 0.20533a_1^{1/2}a_2^{1/2}(a_2-a_1)^{-1}c&<14.5785c<3.645b^5\\
 0.01685a_1(a_1')^{-1}a_2^{-1}(a_2-a_1)^{-2}c&>1.09\cdot 10^6b^{-7}c>272500b^{-3}. 
\end{align*}
Also, $n>0.09226b^{-1/2}c^{1/4}>0.09226\cdot 0.25^{1/4}b^{1/2}$. Applying this to inequalities from Lemmas \ref{lemma: gornja_n_1} and \ref{lema: donja na n u terminima b i c} yields $b\leq 64$ which is a contradiction. 
\end{proof}

\end{document}
